\documentclass{ximera}

% vier voorkeuren die je meteen zelf kan instellen:
\def\uitbr{0} % waarde 1 als je Uitbreiding in de marge wilt, 0 als je dat niet wilt 
\def\wsg{0} % waarde 1 als je verwijzingen naar Wiskunde Samen gevat² in de marge wilt, 0 als je dat niet wilt
\def\rectoverso{0} % waarde 1 als je het PDF-bestand recto-verso wilt laten afdrukken, 0 als je het recto wilt
\def\voetn{1} % waarde 1 als je de voetnoten wilt, 0 als je dat niet wilt 

% marges:
\usepackage[paper=a4paper,margin=3.4cm,marginparwidth=2cm]{geometry}

% packages algemeen:
\pdfOnly{
\usepackage[english,dutch]{babel}
}
\usepackage{amsfonts}
\usepackage{amsmath} %  o.a. \eqref en \text, omgeving equation*
\usepackage{amssymb} % o.a. \nmid
\usepackage{amsthm} % o.a. omgeving proof
\usepackage{graphicx} % o.a. figuren
\usepackage{xcolor} % \color
\usepackage[disable]{todonotes} % \todo 
\usepackage{comment} % omgeving comment
\usepackage{xcolor} % kleuren 
\usepackage[skip=7pt plus1pt, indent=0pt]{parskip} % skip = verticale ruimte tussen twee alinea's, indent = insprong bij nieuwe alinea
\usepackage{multicol} % kolommen
\usepackage{xfrac} % breuken

\usepackage{siunitx} % SI eenheden
\usepackage{eso-pic} % achtergrond voorpagina
\usepackage{tabularx} % kolommen voor rekenmachine
\usepackage{rotating} % voor omgeving turn

% voor figuren met PSTricks:
\usepackage{pstricks} 
\usepackage{pstricks-add}
\usepackage{pst-plot}
\usepackage{pst-node}
\usepackage{pst-coil}
\usepackage{auto-pst-pdf}

% voor figuren met TikZ:
\usepackage{tikz} 
\usepackage{tkz-euclide} 
\usepackage{pgfplots} 
\usetikzlibrary{calc,intersections,through,backgrounds,patterns} 
\pgfplotsset{compat=newest}
\usepgfplotslibrary{fillbetween,colormaps}
\usepackage{import}
\usepackage{tikz-3dplot}

% zelfde parskip in minipage:
\makeatletter
\setlength{\parskip}{\medskipamount}
\newcommand{\@minipagerestore}{\setlength{\parskip}{\medskipamount}}
\makeatother

% voor veranderen van de naam Bibliografie naar Referentielijst:
\pdfOnly{
\addto\captionsdutch{\renewcommand{\bibname}{Referentielijst}}

% voor trefwoordenlijst:
\usepackage{imakeidx}
\makeindex[title=Trefwoordenlijst,program=makeindex,options=-s index.tex,columns=2,intoc=true]
}
% voor nummeren van vergelijkingen:
%WIM%\numberwithin{equation}{chapter} % als je dit desactiveert dan worden vergelijkingen in bijvoorbeeld hoofdstuk 3 genummerd als (1), (2) etc. in plaats van (3.1), (3.2) etc.

% voor hyperlinks en bladwijzers in PDF-bestand:
%WIM%\usepackage[bookmarksopen,bookmarksopenlevel=0,hypertexnames=false,pdfa,bookmarksnumbered]{hyperref} 
%WIM%\usepackage{bookmark} 

% voor hyperlinks: 
\hypersetup{pdfborder={0 0 0}, pdfstartpage=1,linkbordercolor={1 0 0}}%pdfborder={0 0 0} is geen rand rond hyperlinks, pdfborder={1 1 1} wel

% korter commando voor displaystyle, om bijvoorbeeld breuken groter te drukken in doorlopende tekst: 
\newcommand{\D}{\displaystyle}

% (veel)gebruikte verzamelingen:
\newcommand\NN{\mathbb{N}} % verzameling van de natuurlijke getallen
\newcommand\QQ{\mathbb{Q}} % verzameling van de rationale getallen
\newcommand\RR{\mathbb{R}} % verzameling van de re\"ele getallen
\newcommand\ZZ{\mathbb{Z}} % verzameling van de gehele getallen

% (veel)gebruikte operatoren:
\def\co{\operatorname{co}} % coördinaat van een punt
\def\ggd{\operatorname{ggd}} % positieve grootste gemene deler van twee gehele getallen niet beide nul
\def\gr{\operatorname{gr}} % graad van een veelterm

% voor kleur grafieken (donkergroen):
\definecolor{graf}{RGB}{0,100,0} 

% voor lijsten: 
% \usepackage{enumerate}
%WIM%\usepackage[shortlabels]{enumitem} 
\setlist{topsep=0em, itemsep=-0.15em}

% voor small bullet:
\newcommand\sbullet[1][.5]{\mathbin{\vcenter{\hbox{\scalebox{#1}{$\bullet$}}}}}

% voor meer verticale ruimte tussen vergelijkingen in align
\addtolength{\jot}{0.1cm}

% voor meervoudige voetnoten:
\usepackage{fnpct}

% voor het onderdrukken van voetnoten als optie:
\usepackage{letltxmacro}
\LetLtxMacro\Oldfootnote\footnote
\newcommand{\EnableFootNotes}{%
  \LetLtxMacro\footnote\Oldfootnote%
}
\newcommand{\DisableFootNotes}{%
  \renewcommand{\footnote}[2][]{\relax}
}

% voor verticale lijn in de marge (uitbreiding):
\usepackage[framemethod=default]{mdframed}
\usepackage{marginnote}
\reversemarginpar
\ifthenelse{\uitbr < 1}{\definecolor{rood}{RGB}{255,255,255}}{\definecolor{rood}{RGB}{254,64,64}} %HEX: #fe4040
\mdfdefinestyle{uitbreiding}{%
    topline=false,
    rightline=false,
    bottomline=false,
    leftline=true,
    linecolor=rood,
    linewidth=5pt,
    rightmargin=0pt,
    skipabove=10pt,% ipv 3
    skipbelow=0pt,
    leftmargin=-25pt,
    innerleftmargin=20pt,
    innerrightmargin=0pt,
    innertopmargin=0pt,
    innerbottommargin=0pt%,
%	needspace=30pt %minimumhoogte vooraleer lijn wordt gesplitst
    }
\newenvironment{Uitbreiding}{
    \marginpar{
        \center
		\vspace{0.1cm}
		\vspace{7pt}
		\rotatebox{90}{\color{rood}\Large \bf Uitbreiding}
	}
    \begin{mdframed}[style=uitbreiding]
    }{\vspace{-0.05cm}
    \end{mdframed}
}


% omgevingen voor lemma, definitie, voorbeeld etc.:
\newtheoremstyle{mystyle}
    {0em} % Space above
    {0em} % Space below
    {\itshape} % Body font
    {} % Indent amount
    {\bfseries} % Theorem head font
    {.} % Punctuation after theorem head
    {.5em} % Space after theorem head
    {} % Theorem head spec (can be left empty, meaning `normal')
\theoremstyle{mystyle}
%WIM%\newtheorem{lemma}{Lemma}[chapter] % als je [chapter] desactiveert dan Voorbeeld 3 in plaats van Voorbeeld 5.3, als je [chapter] vervangt door [section] dan Voorbeeld 5.2.3 in plaats van Voorbeeld 5.3 
%\theoremstyle{definition} % als je dit activeert dan is wat in de omgeving staat niet cursief gedrukt
\newtheorem{oefening}[lemma]{Oefening}
\newtheorem{definitie}[lemma]{Definitie} 
\newtheorem{voorbeeld}[lemma]{Voorbeeld}
\newtheorem{eigenschap}[lemma]{Eigenschap} 
\newtheorem{stelling}[lemma]{Stelling} 
\newtheorem{gevolg}[lemma]{Gevolg} 
\newtheorem{afspraak}[lemma]{Afspraak} 
\newtheorem{werkwijze}[lemma]{Werkwijze}

% omgeving proof, aangepaste ruimte: 
\makeatletter
\renewenvironment{proof}[1][\proofname]{\par
  \vspace{-\topsep}% remove the space after the theorem
  \pushQED{\qed}%
  \normalfont
  \topsep0pt \partopsep0pt % no space before
  \trivlist
  \item[\hskip\labelsep
        \itshape
    #1\@addpunct{.}]\ignorespaces
}{%
  \popQED\endtrivlist\@endpefalse
  \addvspace{0pt plus 0pt} % no space after
}
\makeatother

% kaderstijlen uit SOHO Wiskunde Plantyn:
\colorlet{steunkleur}{black}
\colorlet{steunkleurlicht}{steunkleur!30!white}
\colorlet{steunkleurkader}{steunkleur!7!white}
\colorlet{steunkleurkaderlicht}{steunkleur!2!white}

\mdfdefinestyle{kaderstijl_vol_licht}
{skipabove=6pt, 
skipbelow=6pt, 
backgroundcolor=steunkleurkader,
linecolor=steunkleurlicht,
linewidth = 0.4pt, 
topline=true,
bottomline=true, 
rightline=true,
innerleftmargin=5pt,
innerrightmargin=5pt,
innertopmargin=5pt,
leftmargin=0cm,
rightmargin=0cm,
innerbottommargin=5pt,
needspace=30pt % minimumhoogte voor splitsen kader
}
\surroundwithmdframed[style=kaderstijl_vol_licht]{definitie}
\surroundwithmdframed[style=kaderstijl_vol_licht]{stelling}
\surroundwithmdframed[style=kaderstijl_vol_licht]{lemma}
\surroundwithmdframed[style=kaderstijl_vol_licht]{eigenschap}
\surroundwithmdframed[style=kaderstijl_vol_licht]{gevolg}

% voor omgevingen voor oefening en antwoord:
\newenvironment{Oefening}{%
    \begin{enumerate}[ 
    series=Oef,
    resume=Oef,
    leftmargin=1.78em,
    label={\bfseries\arabic*.},
    ref=\arabic*
    ]
    \item %
    }{%
    \end{enumerate}
}
\newenvironment{Antwoord}{%
    \begin{enumerate}[%
    series=Antw,
    resume=Antw,
    leftmargin=1.78em,
    label={\bfseries\arabic*.},
    ref=\arabic*
    ]
    \item %
    }{%
    \end{enumerate}
}

% als de omgeving Uitbreiding start met een kaderomgeving (definitie, stelling, eigenschap, lemma of gevolg) dan moet wat extra verticale ruimte voorzien worden, met het commando \uitbreidingstartmetkader:
\def\uitbreidingstartmetkader{\mbox{}\vspace{-0.205cm}}

% % voor schema van de staartdeling:
% \usepackage{stackengine}
% \setstackgap{S}{5pt}
% \stackMath\def\stackalignment{r}
% \newcommand{\myRule}[3][white]{\textcolor{#1}{\rule{#2}{#3}}}
% \let\ph\phantom % enkel voor tekst
% \newcommand{\mph}[1]{% enkel voor math mode
%     \mathcolor{white}{#1}%
% }
% \def\staartmin{\rule{0.25cm}{0.1mm}\myRule{0.3cm}{0.1mm}}
% \def\staartphmin{\myRule{0.25cm}{0.1mm}\myRule{0.3cm}{0.1mm}}
% \newcommand{\staartstreep}[1]{\rule{\widthof{$#1$}}{0.1mm}}
% \newcommand{\staartphstreep}[1]{\myRule{\widthof{$#1$}}{0.1mm}}

% voor kolommen met schema van Horner:
\newcommand{\kolbreed}{1.0cm}
\newcolumntype{H}{>{\centering\arraybackslash$} p{\kolbreed} <{$}}

% voor nieuw commando utikzdashed: onderlijn (zoals underline) maar dan met stippellijn:
\tikzset{
    cheating dash/.code args={on #1 off #2 ends #3}{%
        \csname tikz@addoption\endcsname{%
            \pgfgetpath\currentpath%
            \pgfprocessround{\currentpath}{\currentpath}%
            \csname pgf@decorate@parsesoftpath\endcsname{\currentpath}{\currentpath}%
            \pgfmathparse{max(#1-#3,0)}\let\dashphase=\pgfmathresult%
            \pgfmathparse{\csname pgf@decorate@totalpathlength\endcsname-#1+2*\dashphase}\let\rest=\pgfmathresult%
            \pgfmathparse{#1+#2}\let\onoff=\pgfmathresult%
            \pgfmathparse{max(floor(\rest/\onoff), 1)}\let\nfullonoff=\pgfmathresult%
            \pgfmathparse{max((\rest-\onoff*\nfullonoff)/\nfullonoff+#2, #2)}\let\offexpand=\pgfmathresult%
            \pgfsetdash{{#1}{\offexpand}}{\dashphase pt}}%
    },
    cheating dash per segment/.style args={on #1 off #2 ends #3}{
        /utils/exec=\csname tikz@options\endcsname,%
        decoration={show path construction,
            lineto code={\draw [cheating dash=on #1 off #2 ends #3] (\tikzinputsegmentfirst) -- (\tikzinputsegmentlast);},
            curveto code={\draw [cheating dash=on #1 off #2 ends #3] (\tikzinputsegmentfirst) .. controls (\tikzinputsegmentsupporta) and (\tikzinputsegmentsupportb) .. (\tikzinputsegmentlast);},
            closepath code={\draw [cheating dash=on #1 off #2 ends #3] (\tikzinputsegmentfirst) -- (\tikzinputsegmentlast);}
        },
        decorate,
    },
}
\newcommand{\utikzdash}[1]{%
    \tikz[baseline=(todotted.base)]{
        \node[inner sep=0pt,outer sep=1.5pt] (todotted) {#1};
        \draw[cheating dash per segment=on 2pt off 2pt ends 2pt, line width=0.4pt] (todotted.south west) -- (todotted.south east);
    }%
}%

\providecommand{\xmemph}[1]{\textit{#1}}

% voor nieuw commando underdashed: onderlijn met stippellijnen en ook het woord cursief zetten:
\newcommand{\underdashed}[1]{%
    {\em\utikzdash{\!#1\!}}%
}

% voor icoon TI-84 Plus met verwijzing naar filmpje:
\newcommand{\grmlink}{\raisebox{0cm}{\includegraphics[width=1cm]{TI84Plus-icoon}}}
\newcommand{\grmref}[1]{%
	\reversemarginpar%
	\marginpar{
		\vspace{-0.4cm}%
		\htmladdnormallink{\grmlink}{#1}
	}	
}

% GRM knoppen:
\newcommand{\GRM}[1]{\fbox{\rule[0mm]{0cm}{0.215cm}\textup{\texttt{#1}}}}
\newcommand{\wedgetext}{{\raisebox{0.02cm}{\begin{turn}{90}>\end{turn}}}}
\newcommand{\veetext}{\raisebox{0.2cm}{\begin{turn}{-90}>\end{turn}}}

% voor kolommen met GRM screens:
\newlength{\widthallscreens}
\newlength{\widthscreens}
\newlength{\spaceleftscreen}
\newlength{\spacerightscreen}
\newlength{\spacebetweenscreens}

\newcommand{\setscreens}{
	\setlength{\spaceleftscreen}{2pt}
	\setlength{\spacerightscreen}{2pt}
	\setlength{\spacebetweenscreens}{8pt}
	\addtolength{\linewidth}{-28pt} % 3*8pt tussen vier screens en 2 pt links en 2 pt rechts
	\setlength{\widthallscreens}{\linewidth}
	\setlength{\widthscreens}{0.25\widthallscreens}
	\addtolength{\linewidth}{28pt}
	\newcolumntype{G}{p{\widthscreens}}
	\newcolumntype{s}{p{\spacebetweenscreens}}
	\newcolumntype{L}{p{\spaceleftscreen}}
	\newcolumntype{R}{p{\spacerightscreen}}
	\setlength{\tabcolsep}{0pt}
}

% voor icoon met verwijzing naar Wiskunde Samen gevat² in de marge:
% \newcommand{\wsglink}{\raisebox{0cm}{\includegraphics[width=1cm]{wsglogo}}}
\newcommand{\wsglink}{\raisebox{0cm}{LOGO}}
\newcommand{\htmladdnormallink}[2]{\href{#2}{#1}}
\newcounter{pagnrwsg}
\newcommand{\wsgref}[3]{% 
    % #1 woord dat onderstippeld wordt
    % #2 pagina van Wiskunde Samen gevat² waar je op terecht komt als je op de link klikt  
    % #3 afgedrukt op het logo van Wiskunde Samen gevat² in de marge: één of meerdere paginanummers
    \ifthenelse{\wsg < 1}{#1}{%
        \underdashed{#1}%
        {\setcounter{pagnrwsg}{#2}}%
        {\addtocounter{pagnrwsg}{14}}%
        \reversemarginpar%
        \marginpar{\vspace{-0.6cm}%
           \htmladdnormallink{\wsglink}{https://online.fliphtml5.com/sanky/laea/\#p=\arabic{pagnrwsg}}%
            \raisebox{0.41cm}[0cm][0cm]{%
                \hspace{-1.5cm}\makebox[2cm][c]{\colorbox{white}{\texttt{\footnotesize{#3}}}}%
            }%
        }%
    }%
}

% voor invoegen blanco pagina bij de optie recto-verso:
\def\blancobijrectoverso{
	\ifthenelse{\rectoverso < 1}{\clearpage}{
	\clearpage
	\thispagestyle{empty}
	\mbox{}
	\clearpage
	}
}

% aantal pagina's van het bestand:
\ifthenelse{\rectoverso < 1}{\def\totpag{57}}{\def\totpag{66}}

% documenteigenschappen:
% \usepackage{hyperxmp}
% \hypersetup{
% pdftitle={Open Source Wiskunde Aan zet: Veeltermen}, 
% pdfnumpages={\totpag},
% pdfauthor={Koen De Naeghel},
% pdflang={nl},
% pdfkeywords={wiskunde, open source, wiskunde aan zet, veeltermen, secundair onderwijs, tweede graad, doorstroomfinaliteit},
% pdfsubject={Veeltermen},
% pdfcopyright={\unichar{"24B8} 2024 Koen De Naeghel},
% pdfdate={17 december 2024},
% pdfapart=1,
% pdfstartview=}

% voor het benoemen van titel, auteur en datum:
% \title{Veeltermen}
% \author{Auteur: Koen De Naeghel}
% \date{\today}

\newcommand\BackgroundPic{
    \put(-260,-125){
    \parbox[b][\paperheight]{\paperwidth}{%
    \vfill
    \centering
    \includegraphics[height=\paperwidth, keepaspectratio]{WaZlogo}%
    \vfill
}}}

\addPrintStyle{..}

\begin{document}
	\author{Wim Obbels}
	\xmtitle{Vierkantswortels van complexe getallen}{}
	\label{xim:cmplx_vierkantswortels}

    Een belangrijke motivatie voor het invoeren van complexe getallen was het creëren van een 'vierkantswortel uit $-1$', die we $i$ noteerden.
    Het is nu tijd om niet alleen de vierkantswortel van $-1$ te bestuderen, maar die van willekeurige complexe getallen.


    \begin{expandable}{remark}{Vierkantswortels van reële getallen}
        Je bent al lang vertrouwd met vierkantswortels van reële getallen. We herhalen hier kort de belangrijkste eigenschappen.

        \begin{definition}[Vierkantswortel van een reëel getal]\nl

            Een reëel getal $w$ is een {vierkantswortel} van een reëel getal $a$ als 
            \formulevb{a = w^2}{2 \text{ en } -2 \text{ zijn wortels van} 4 \text{ want } 2^2 = 4 \text{ en } (-2)^2 = 4}
        \end{definition}

        Enkel positieve getallen hebben vierkantswortels, en wel steeds twee die aan elkaar tegengesteld zijn. Bij $0$ spreken we af dat de twee tegengestelde wortel samenvallen.
        \begin{proposition}\nl

            \begin{tabular}{@{}l@{ }l}
            Een reëel getal $a$ heeft een vierkantswortel & als en slechts als $a$ positief is (dus $0\leq a$) \\
                                                          & als en slechts als vergelijking $x^2 = a$ een oplossing heeft \\
                                                          & als en slechts als vergelijking $x^2 - a = 0 $ een oplossing heeft 
            \end{tabular}

            Als $a>0$, dan heeft $a$ twee vierkantswortels, een positieve die we noteren met $\sqrt{a}$, en een negatieve die we noteren met $-\sqrt{a}$.
        \end{proposition}

        \begin{quickquestion*}
            Geeft telkens de precieze verzameling $x$'en waarvoor de gelijkheid wel geldt.
            \begin{itemize}
                \item Geldt voor elke reëel getal $x$ dat $(\sqrt{x})^2 = x$?
                \item Geldt voor elke reëel getal $x$ dat $(\sqrt{|x|})^2 = x$?
                \item Geldt voor elke reëel getal $x$ dat $(\sqrt{|x|})^2 = |x|$?
                \item Geldt voor elke reëel getal $x$ dat $|\sqrt{|x|}|^2 = |x|$?
                \item 
            \end{itemize}
         \end{quickquestion*}


        Er is een evident verband met grafieken van functies: de vergelijking $x^2-a=0$ is oplosbaar als de grafiek van de functie $y=x^2-a$ de $x$-as snijdt, en de nulwaarden of nulpunten zijn precies de twee vierkantswortels van $a$.

    \end{expandable}


    De definitie van vierkantswortels van complexe getallen is volledig gelijkaardig aan de gekende definitie voor reële getallen:
    \begin{definition}[Vierkantswortel van complex getal]\nl

        Een complex getal $w$ is een {vierkantswortel} van een complex getal $z$ als 
        \formulevb{z = w^2}{i \text{ en } -i \text{ zijn wortels van} -1 \text{ want } i^2 = (-i)^2 = -1}
    \end{definition}

    \begin{example}
        \begin{question} Het complex getal $2i$ is \choiceEen vierkantswortel van $-4$. \end{question}
        \begin{question} Het complex getal $-2i$ is \choiceEen vierkantswortel van $-4$. \end{question}
        \begin{question} Het complex getal $2$ is \choiceGeen vierkantswortel van $4$. \end{question}
        \begin{question} Het complex getal $-2$ is \choiceGeen vierkantswortel van $-4$. \end{question}
        \begin{question} Het complex getal $1+2i$ is \choiceEen vierkantswortel van $-1+4i$. \end{question}
        \begin{question} Het complex getal $1-2i$ is \choiceGeen vierkantswortel van $-1+4i$. \end{question}
        \begin{question} Het complex getal $1+i$ is \choiceEen vierkantswortel van $2i$. \end{question}
    \end{example}

    Bij complexe getallen zijn vierkantswortels eenvoudiger dan bij reële getallen, want het verschil tussen positieve en negatieve getallen valt weg.

    \begin{proposition}
        Elke complex getal $w$ heeft twee tegengestelde vierkantswortels.

        of ook, equivalent geformuleerd:
        {\centering

        De vergelijking $x^2-w = 0$ heeft steeds twee aan elkaar tekengestelde wortels.
        }
    \end{proposition}
    We spreken af dat voor $w=0$ het getal $0$ een 'dubbele wortel' is, waarbij $0$ en $-0=0$ tegengestelden zijn\ldots

    \begin{expandable}{proof}{Een erg gedetailleerd bewijs}
        We zullen de eigenschap bewijzen door voor een willekeurig complex getal $w=a+bi$ expliciet twee vierkantswortels te berekenen.

        We zoeken twee oplossingen van de vergelijking $z^2-w=0$. Stel $z=x+iy$, dan zoeken we dus geschikte reële getallen $x$ en $y$ zodat 
        $$
        a+bi = (x+iy)^2 = x^2+2xyi+(yi)^2 = x^2-y^2 +2xyi.
        $$
        Omdat twee complexe getallen enkel aan elkaar gelijk zijn als hun reële delen en imaginaire delen gelijk zijn, volgt dat we oplossingen zoeken van het stelsel
        $$
        (A) \leftrightarrow \begin{cases} x^2 - y^2 = a & (1) \\
                        2xy = b & (2)
        \end{cases}
        $$
        Dit is een stelsel van twee vergelijkingen met twee onbekenden, maar de vergelijkingen zijn niet lineair. De theorie van het oplossen van lineaire stelsels helpt ons dus niet.

        Het oplossen van niet-lineaire stelsels is in het algemeen hopeloos, maar hier is er een trucje.
        Bij gebrek aan betere ideeën, berekenen we $a^2+b^2$ door vergelijkingen (1) en (2) te substitueren
        $$
        a^2+b^2 = (x^2-y^2)^2 + (2xy)^2 = x^4 -2x^2y^2 + y^4 + 4 x^2y^2 = x^4+2x^2y^2 + y^4 = (x^2+y^2)^2
        $$
        Dat is een mooi resultaat, maar helpt het ons ook vooruit?

        We kunnen nu vergelijking $x^2+y^2=\sqrt{a^2+b^2}$ toevoegen aan het stelsel (A) zonder de oplossingen van dat stelsel te wijzigen.

        \begin{quickquestion*}{} Waarom kunnen er geen oplossingen verdwijnen, en ook geen oplossingen bijkomen.\end{quickquestion*}

        We zoeken dus $x$ en $y$ zodat
        \begin{equation}
        \begin{cases} x^2 - y^2 = a & (1) \\
                        2xy = b & (2) \\
                     x^2 + y^2 = \sqrt{a^2+b^2} & (3)
        \end{cases} 
        \end{equation}
        Op het eerste zicht hebben we het probleem moeilijke gemaakt, want nu hebben we zelfs een vergelijking meer om mee rekening te hopuden.
        Maar door de vergelijkingen (1) en (3) bij elkaar op te tellen resp. van elkaar af te trekken, vinden we relatief eenvoudig uitdrukkingen voor $x^2$ en $y^2$:
        $$
        \begin{cases} 2x^2 = a + \sqrt{a^2+b^2} \\
                    -2y^2 = a -\sqrt{a^2+b^2} 
        \end{cases}
        , \text{ of dus } (B) \leftrightarrow 
        \begin{cases} x = \pm\sqrt{\frac{a + \sqrt{a^2+b^2}}{2}} \\
                      y = \pm\sqrt{\frac{-a +\sqrt{a^2+b^2}}{2}} 
        \end{cases}
        $$

        Je kan opmerken dat deze uitdrukkingen zinvol zijn, omdat elke uitdrukking onder een wortelteken steeds zeker positief is.
        Bovendien voldoen de $x$ en $y$ uit de laatste uitdrukkingen bijna automatisch aan de vergelijking $2xy=b$. 
        Dat kunnen we makkelijkst nagaan door $4x^2y^2$ te berekenen (want dan hebben we minder wortels):
        $$
        4x^2y^2 = (a+\sqrt{a^2+b^2}(-a+\sqrt{a^2+b^2})) = -a^2 + (a^2+b^2) = b^2
        $$
        en dus is $2xy=\pm b$. 
        Er zijn in (B) nog twee mogelijkheden voor $x$, en twee voor $y$. Door deze slim te combineren kunne we er voor zorgen dat inderdaad $2xy=b$ (en niet $-b$).
        
        Conclusie: we vinden voor elk complex getal $w=a+bi$ twee vierkantswortels $z=\pm(x+iy)$, namelijk

        Als $b\geq 0$ dan zijn $z_1$ en $-z_1$ wortels met $z_1 = \sqrt{\frac{a + \sqrt{a^2+b^2}}{2}} + i\sqrt{\frac{-a + \sqrt{a^2+b^2}}{2}}$, en 

        als $b< 0$ dan zijn $z_2$ en $-z_2$ wortels met $z_2 = \sqrt{\frac{a + \sqrt{a^2+b^2}}{2}} - i\sqrt{\frac{-a + \sqrt{a^2+b^2}}{2}}$, en 

        Dit beëindigt het bewijs: we hebben voor elk complex getal $w$ expliciet twee wortels berekend.
    \end{expandable}

    \begin{remark}
        Vierkantswortels in $\C$ zijn eigenlijk eenvoudiger dan in $\R$, omdat er \textit{altijd} vierkantswortels bestaan.
        Toch is er een subtiliteit.

        Inderdaad, voor ëen positief reëel getal $a$ hebben we $\sqrt{a}$ gedefinieerd als de \textit{positieve} vierkantswortel, en $-\sqrt{a}$ is dan natuurlijk de negatieve wortel.
        Deze keuze is voor complexe getallen echter onmogelijk, omdat op $\C$ de begrippen 'positief' en 'negatief' geen betekenis hebben.

        Het is dus niet eenvoudig mogelijk om uit de twee wortels van een complex getal $w$ er op consistente wijze één te kiezen om $\sqrt{w}$ te zijn.
        \\
        Naast het probleem met de \hyperref[eig:geen_wortelteken_voor_complexe_getallen]{rekenregels}, is dit een extra reden om \textbf{het gebruik van het symbool $\sqrt{\cdot}$ voor complexe getallen te vermijden}.

    \end{remark}
\end{document}